\documentclass[11pt,a4paper]{article}

\usepackage[utf8]{inputenc}
\usepackage[T1]{fontenc}
\usepackage{amsmath,amssymb,amsthm}
\usepackage{geometry}
\usepackage{enumitem}
\usepackage{hyperref}

\geometry{margin=1in}

\theoremstyle{plain}
\newtheorem{theorem}{Theorem}

\title{Problem Set 4}
\author{}
\date{}

\begin{document}

\maketitle

\begin{enumerate}[leftmargin=*]

\item Quantum antimatter fuel comes in small pellets, which is convenient since the many moving parts of the LAMBCHOP each need to be fed fuel one pellet at a time. However, minions dump pellets in bulk into the fuel intake. You need to figure out the most efficient way to sort and shift the pellets down to a single pellet at a time. The fuel control mechanisms have three operations:
\begin{itemize}
    \item Add 1 fuel pellet
    \item Remove 1 fuel pellet
    \item Divide the entire group of fuel pellets by 2 (due to the destructive energy released when a quantum antimatter pellet is cut in half, the safety controls will only allow this to happen if there is an even number of pellets)
\end{itemize}

Write a function called $\texttt{answer}(n)$ which takes a positive integer $n$ as a string and returns the minimum number of operations needed to transform the number of pellets to 1. Prove the correctness of your algorithm. (E.g. $29 \to 28 \to 14 \to 7 \to 8 \to 4 \to 2 \to 1$)

\item (Exercise 4.3 from the textbook) You are consulting for a trucking company that does a large amount of business shipping packages between New York and Boston. The volume is high enough that they have to send a number of trucks each day between the two locations. Trucks have a fixed limit $W$ on the maximum amount of weight they are allowed to carry. Boxes arrive at the New York station one by one, and each package $i$ has a weight $w_i$. The trucking station is quite small, so at most one truck can be at the station at any time. Company policy requires that boxes are shipped in the order they arrive; otherwise, a customer might get upset upon seeing a box that arrived after his make it to Boston faster. At the moment, the company is using a simple greedy algorithm for packing; they pack boxes in the order they arrive, and whenever the next box does not fit, they send the truck on its way.

But they wonder if they might be using too many trucks, and they want your opinion on whether the situation can be improved. Here is how they are thinking. Maybe one could decrease the number of trucks needed by sometimes sending off a truck that was less full, and in this way allow the next few trucks to be better packed.

Prove that, for a given set of boxes with specified weights, the greedy algorithm currently in use actually minimizes the number of trucks that are needed. Your proof should follow the type of analysis we used for the Interval Scheduling Problem: it should establish the optimality of this greedy packing algorithm by identifying a measure under which it ``stays ahead'' of all other solutions.

\end{enumerate}

\end{document}